% GENERATED FILE. Edit canonical lessons, then run npm run generate:books.
% canonical-source-hash: fnv1a64:d010aca70fc59482
% canonical-lessons: IT-C312-antenato, IT-C312-posteri, IT-C312-discendenza, IT-C312-stirpe, IT-C312-albero-genealogico, IT-R312-ripasso-305-312

\chapter{Le radici}
\label{ch:it-le-radici}
\hlchaptermodality{312}

\begin{chapteropening}
\textbf{By the end of this chapter:} \emph{I can talk about where a family comes from: ancestors, those who come after, the line of descent and the family tree.}

Complete the last lesson of chapter 312: I can talk about where a family comes from: ancestors, those who come after, the line of descent and the family tree.
\end{chapteropening}

\section[l'antenato]{l'antenato — the ancestor}

\label{lesson:IT-C312-antenato}



\begin{quote}
Before the new one: say the Italian for the bouquet, then the Italian for the wedding favour.
\end{quote}

\subsection*{You'll want to know: l'antenato}
\textbf{l'antenato} — \textquotedblleft{}the ancestor\textquotedblright{}.

\begin{quote}
I miei antenati erano contadini.
\end{quote}

\begin{quote}
\emph{My ancestors were farmers.}
\end{quote}

Plural \emph{gli antenati}.

\begin{cousinweb}
\emph{antenato} is usually traced to Late Latin \emph{antenātus} \textquotedblleft{}born before\textquotedblright{}, a learned word built on \emph{ante} \textquotedblleft{}before\textquotedblright{} and \emph{nātus} \textquotedblleft{}born\textquotedblright{} (Lewis \& Short).
\end{cousinweb}

\begin{grammarlens}[title={What you've built}]
One more word for telling, explaining and describing at B1.
\end{grammarlens}

\subsection*{Your turn}
\begin{itemize}
\raggedright
  \item \emph{Say it:} \emph{l'antenato}
  \item \emph{Say it:} \emph{l'antenato}, once more
  \item \emph{Say it:} a sentence of your own with \emph{l'antenato}
  \item \emph{Recall:} say the Italian for the bouquet, then the Italian for the wedding favour, then say \emph{l'antenato} again
\end{itemize}

\begin{tcolorbox}[breakable,colback=teal!4,colframe=teal!35!black,title={Before you move on}]
What does l'antenato mean? (\textquotedblleft{}The ancestor\textquotedblright{}.) Say it once more.
\end{tcolorbox}
\section[i posteri]{i posteri — posterity, future generations}

\label{lesson:IT-C312-posteri}



\begin{quote}
Before the new one: say the Italian for the wedding favour, then the Italian for the ancestor.
\end{quote}

\subsection*{You'll want to know: i posteri}
\textbf{i posteri} — \textquotedblleft{}posterity, future generations\textquotedblright{}.

\begin{quote}
Questa casa resterà ai posteri.
\end{quote}

\begin{quote}
\emph{This house will be left to posterity.}
\end{quote}

Plural only; formal.

\begin{cousinweb}
\emph{posteri} is a learned borrowing from Latin \emph{posterī} \textquotedblleft{}those coming after\textquotedblright{}, from \emph{posterus} \textquotedblleft{}following\textquotedblright{} (Lewis \& Short).
\end{cousinweb}

\begin{grammarlens}[title={What you've built}]
One more word for telling, explaining and describing at B1.
\end{grammarlens}

\subsection*{Your turn}
\begin{itemize}
\raggedright
  \item \emph{Say it:} \emph{i posteri}
  \item \emph{Say it:} \emph{i posteri}, once more
  \item \emph{Say it:} a sentence of your own with \emph{i posteri}
  \item \emph{Recall:} say the Italian for the wedding favour, then the Italian for the ancestor, then say \emph{i posteri} again
\end{itemize}

\begin{tcolorbox}[breakable,colback=teal!4,colframe=teal!35!black,title={Before you move on}]
What does i posteri mean? (\textquotedblleft{}Posterity, future generations\textquotedblright{}.) Say it once more.
\end{tcolorbox}
\section[la discendenza]{la discendenza — the descendants, the line of descent}

\label{lesson:IT-C312-discendenza}



\begin{quote}
Before the new one: say the Italian for the ancestor, then the Italian for posterity.
\end{quote}

\subsection*{You'll want to know: la discendenza}
\textbf{la discendenza} — \textquotedblleft{}the descendants, the line of descent\textquotedblright{}.

\begin{quote}
Ha una discendenza numerosa: dieci nipoti.
\end{quote}

\begin{quote}
\emph{He has many descendants: ten grandchildren.}
\end{quote}

Built on \emph{discendere}, ‘to go down’.

\begin{cousinweb}
\emph{discendenza} is built on \emph{discendere} \textquotedblleft{}to descend\textquotedblright{}, from Latin \emph{dēscendere} (EtymDB), from \emph{dē} and \emph{scandere} \textquotedblleft{}to climb\textquotedblright{}.
\end{cousinweb}

\begin{grammarlens}[title={What you've built}]
One more word for telling, explaining and describing at B1.
\end{grammarlens}

\subsection*{Your turn}
\begin{itemize}
\raggedright
  \item \emph{Say it:} \emph{la discendenza}
  \item \emph{Say it:} \emph{la discendenza}, once more
  \item \emph{Say it:} a sentence of your own with \emph{la discendenza}
  \item \emph{Recall:} say the Italian for the ancestor, then the Italian for posterity, then say \emph{la discendenza} again
\end{itemize}

\begin{tcolorbox}[breakable,colback=teal!4,colframe=teal!35!black,title={Before you move on}]
What does la discendenza mean? (\textquotedblleft{}The descendants, the line of descent\textquotedblright{}.) Say it once more.
\end{tcolorbox}
\section[la stirpe]{la stirpe — the lineage, the stock}

\label{lesson:IT-C312-stirpe}



\begin{quote}
Before the new one: say the Italian for posterity, then the Italian for the descendants.
\end{quote}

\subsection*{You'll want to know: la stirpe}
\textbf{la stirpe} — \textquotedblleft{}the lineage, the stock\textquotedblright{}.

\begin{quote}
Viene da una stirpe di pescatori.
\end{quote}

\begin{quote}
\emph{He comes from a line of fishermen.}
\end{quote}

Literary.

\begin{cousinweb}
\emph{stirpe} is a learned borrowing from Latin \emph{stirps} \textquotedblleft{}a stem, a stock, a family\textquotedblright{} (Lewis \& Short).
\end{cousinweb}

\begin{grammarlens}[title={What you've built}]
One more word for telling, explaining and describing at B1.
\end{grammarlens}

\subsection*{Your turn}
\begin{itemize}
\raggedright
  \item \emph{Say it:} \emph{la stirpe}
  \item \emph{Say it:} \emph{la stirpe}, once more
  \item \emph{Say it:} a sentence of your own with \emph{la stirpe}
  \item \emph{Recall:} say the Italian for posterity, then the Italian for the descendants, then say \emph{la stirpe} again
\end{itemize}

\begin{tcolorbox}[breakable,colback=teal!4,colframe=teal!35!black,title={Before you move on}]
What does la stirpe mean? (\textquotedblleft{}The lineage, the stock\textquotedblright{}.) Say it once more.
\end{tcolorbox}
\section[l'albero genealogico]{l'albero genealogico — the family tree}

\label{lesson:IT-C312-albero-genealogico}



\begin{quote}
Before the new one: say the Italian for the descendants, then the Italian for the lineage.
\end{quote}

\subsection*{You'll want to know: l'albero genealogico}
\textbf{l'albero genealogico} — \textquotedblleft{}the family tree\textquotedblright{}.

\begin{quote}
Mio padre ha disegnato l'albero genealogico della famiglia.
\end{quote}

\begin{quote}
\emph{My father drew the family tree.}
\end{quote}

\emph{Genealogico} — of the family line.

\begin{cousinweb}
\emph{albero} is inherited from Latin \emph{arbor} \textquotedblleft{}tree\textquotedblright{} (EtymDB); \emph{genealogico} comes from Greek \emph{genealogía}, an account of a family line, which Latin took over as \emph{geneālogia} (Lewis \& Short).
\end{cousinweb}

\begin{grammarlens}[title={What you've built}]
One more word for telling, explaining and describing at B1.
\end{grammarlens}

\subsection*{Your turn}
\begin{itemize}
\raggedright
  \item \emph{Say it:} \emph{l'albero genealogico}
  \item \emph{Say it:} \emph{l'albero genealogico}, once more
  \item \emph{Say it:} a sentence of your own with \emph{l'albero genealogico}
  \item \emph{Recall:} say the Italian for the descendants, then the Italian for the lineage, then say \emph{l'albero genealogico} again
\end{itemize}

\begin{tcolorbox}[breakable,colback=teal!4,colframe=teal!35!black,title={Before you move on}]
What does l'albero genealogico mean? (\textquotedblleft{}The family tree\textquotedblright{}.) Say it once more.
\end{tcolorbox}
\section[(dialogue)]{Ripasso dei capitoli 305–312}

\label{lesson:IT-R312-ripasso-305-312}



\begin{quote}
Say the Italian for the flood and to flood.
\end{quote}

\subsection*{Your turn}
Say these aloud:

\begin{itemize}
\raggedright
  \item the flood, to flood, the landslide, the avalanche, to drown
  \item the rescue, the lifeguard, the shipwreck, missing, to survive
  \item the birth, childbirth, the first-born child, to wean, the dummy
  \item the christening, to christen, the godfather, the godmother, the confirmation
  \item youth, of age, underage, the generation, to hand down
  \item to live together, to move, to fit in, to settle in, to return to one's country
  \item to get married, the wedding ring, the veil, the bouquet, the wedding favour
  \item the ancestor, posterity, the descendants, the lineage, the family tree
\end{itemize}

\begin{tcolorbox}[breakable,colback=teal!4,colframe=teal!35!black,title={Before you move on}]
The flood? (\textbf{l'alluvione}.) To flood? (\textbf{allagare}.) The landslide? (\textbf{la frana}.) The avalanche? (\textbf{la valanga}.) To drown? (\textbf{annegare}.) The rescue? (\textbf{il salvataggio}.) The lifeguard? (\textbf{il bagnino}.) The shipwreck? (\textbf{il naufragio}.) Missing? (\textbf{disperso}.) To survive? (\textbf{sopravvivere}.) The birth? (\textbf{la nascita}.) Childbirth? (\textbf{il parto}.) The first-born child? (\textbf{il primogenito}.) To wean? (\textbf{svezzare}.) The dummy? (\textbf{il ciuccio}.) The christening? (\textbf{il battesimo}.) To christen? (\textbf{battezzare}.) The godfather? (\textbf{il padrino}.) The godmother? (\textbf{la madrina}.) The confirmation? (\textbf{la cresima}.) Youth? (\textbf{la giovinezza}.) Of age? (\textbf{maggiorenne}.) Underage? (\textbf{minorenne}.) The generation? (\textbf{la generazione}.) To hand down? (\textbf{tramandare}.) To live together? (\textbf{convivere}.) To move? (\textbf{trasferirsi}.) To fit in? (\textbf{integrarsi}.) To settle in? (\textbf{ambientarsi}.) To return to one's country? (\textbf{rimpatriare}.) To get married? (\textbf{sposarsi}.) The wedding ring? (\textbf{la fede}.) The veil? (\textbf{il velo}.) The bouquet? (\textbf{il bouquet}.) The wedding favour? (\textbf{la bomboniera}.) The ancestor? (\textbf{l'antenato}.) Posterity? (\textbf{i posteri}.) The descendants? (\textbf{la discendenza}.) The lineage? (\textbf{la stirpe}.) The family tree? (\textbf{l'albero genealogico}.)
\end{tcolorbox}
